\documentclass[11pt]{article}

\usepackage[margin=1in]{geometry}
\usepackage[T1]{fontenc}
\usepackage[utf8]{inputenc}
\usepackage{lmodern}
\usepackage{microtype}
\usepackage{amsmath,amssymb,amsthm,mathtools}
\usepackage[most]{empheq}
\usepackage{booktabs}
\usepackage{array}
\usepackage{longtable}
\usepackage{enumitem}
\usepackage{xcolor}
\usepackage{framed}
\usepackage{url}
\usepackage{hyperref}
\urlstyle{tt}
% \fp{} typesets a file path in typewriter, breakable at / and _ , no escaping needed
\newcommand{\fp}{\begingroup\urlstyle{tt}\Url}
\setlength{\emergencystretch}{3em}

\allowdisplaybreaks
\setlength{\parindent}{0pt}
\setlength{\parskip}{0.5em}

%%%%%%%%%%%%%%%%%%%%%%%%%%%%%%%%%%%%%%%%%%%%%%%%%%%%%%%%%%%%%%%%%%%%%%%%%%%%%%%
% Macros
%%%%%%%%%%%%%%%%%%%%%%%%%%%%%%%%%%%%%%%%%%%%%%%%%%%%%%%%%%%%%%%%%%%%%%%%%%%%%%%
\newcommand{\dd}{{\rm d}}
\newcommand{\ee}{{\rm e}}
\newcommand{\E}{\mathbb{E}}
\newcommand{\Prob}[1]{\mathbb{P}\!\left\{#1\right\}}
\newcommand{\xt}{X_t}
\newcommand{\wt}{W_t}
\newcommand{\Ihat}{\widehat{I}}
\newcommand{\Kb}{\mathcal{K}}
\newcommand{\Icell}{\mathcal{I}}
\newcommand{\Vfree}{\mathcal{V}}
\newcommand{\Lap}[1]{\widetilde{#1}}
\newcommand{\ddt}[1]{\frac{\dd #1}{\dd t}}
\newcommand{\Fhyp}{{}_2F_1}
\newcommand{\Aa}{A}
\newcommand{\Bb}{B}

\theoremstyle{plain}
\newtheorem{theorem}{Result}[section]
\newtheorem{proposition}[theorem]{Result}
\newtheorem{lemma}[theorem]{Lemma}
\newtheorem{corollary}[theorem]{Corollary}
\theoremstyle{definition}
\newtheorem{definition}[theorem]{Definition}
\theoremstyle{remark}
\newtheorem{remark}[theorem]{Remark}
\newtheorem{task}{Task}

% Errata tags
\newcommand{\ERR}{\textcolor{red!70!black}{\textbf{[ERROR]}}}
\newcommand{\GAP}{\textcolor{orange!80!black}{\textbf{[GAP]}}}
\newcommand{\NOT}{\textcolor{blue!60!black}{\textbf{[NOTATION]}}}
\newcommand{\TYP}{\textcolor{gray!80!black}{\textbf{[TYPO]}}}
\newcommand{\STR}{\textcolor{purple!70!black}{\textbf{[STRUCTURE]}}}

\hypersetup{
  colorlinks=true,
  linkcolor=blue!55!black,
  citecolor=red!55!black,
  urlcolor=blue!60!black,
  pdftitle={Working notes: a burst-aware extension of the Basic Model of Viral Replication},
  pdfauthor={Working notes}
}

\title{\textbf{Working notes}\\[0.4em]
A burst-aware extension of the Basic Model of Viral Replication\\[0.3em]
\large Consolidated BDC results, the renewal reformulation,\\
and an errata catalogue for Chapters 1, 3, 4, 5 and 6}
\author{}
\date{Compiled \today}

\begin{document}
\maketitle
\vspace{-3em}

\tableofcontents
\clearpage

%%%%%%%%%%%%%%%%%%%%%%%%%%%%%%%%%%%%%%%%%%%%%%%%%%%%%%%%%%%%%%%%%%%%%%%%%%%%%%%
\section{Purpose and scope}
%%%%%%%%%%%%%%%%%%%%%%%%%%%%%%%%%%%%%%%%%%%%%%%%%%%%%%%%%%%%%%%%%%%%%%%%%%%%%%%

These notes do three things.

\begin{enumerate}[leftmargin=1.6em]
\item \textbf{Part I} consolidates the single-cell birth--death--catastrophe (BDC) results
scattered across Chapters 3 and 4 into one verified, internally consistent set. Several
formulas are simplified, several are corrected, and several results that the chapters
have all the ingredients for but never state are stated here.

\item \textbf{Part II} sets out the extension of the Basic Model of Viral Replication (BMVR)
that Chapter~4 \S9 was reaching for. The attempt in \S9 as written cannot be repaired ---
it is dimensionally inconsistent --- but the object it was groping towards exists,
is exactly solvable, and yields several results of publication value.

\item \textbf{Part IV} is a systematic catalogue of every error, gap, notational collision
and structural problem found in the documents read. It is intended to be worked through
and ticked off.
\end{enumerate}

Everything asserted in Parts I and II has been checked, either algebraically or by
independent numerical evaluation at $(\beta,\mu,\delta)=(1,0.2,0.05)$, $(1,0,0.1)$,
$(1,0.9,0.1)$ and $(1,0.5,1/3)$. Where a result is \emph{not} yet established, it is
flagged as a Task in Part III rather than stated as a result.

%%%%%%%%%%%%%%%%%%%%%%%%%%%%%%%%%%%%%%%%%%%%%%%%%%%%%%%%%%%%%%%%%%%%%%%%%%%%%%%
\section{Canonical notation for these notes}
%%%%%%%%%%%%%%%%%%%%%%%%%%%%%%%%%%%%%%%%%%%%%%%%%%%%%%%%%%%%%%%%%%%%%%%%%%%%%%%

The thesis currently uses $\lambda$ and $\beta$ interchangeably for the birth rate,
$\delta$ for both the catastrophe rate and the infected-cell death rate, $K$ for both the
second moment and the burst size, $H$ and $R$ for the same absorbing state, and $I$ for both
a probability and a cell count. These collisions become fatal the moment the intracellular
and population models are joined, which is exactly what Part II does. The following
choices are used throughout and should be adopted in the thesis.

\begin{center}
\begin{tabular}{@{}llp{0.42\textwidth}@{}}
\toprule
Symbol & Meaning & Thesis usage \\
\midrule
$\beta$ & intracellular birth rate (per capita) & $\lambda$ in Ch.3 \S3, $\beta$ in Ch.3 \S4--\S11 \\
$\mu$ & intracellular death rate (per capita) & same \\
$\delta$ & catastrophe (burst) rate (per capita) & same; but also infected-cell death in Ch.1 \\
$\xt$ & intracellular count & same \\
$R$ & absorbing rupture state & $H$ in Ch.3, $R$ in Ch.4/Ch.5 \\
$\wt$ & released (extracellular) count from one cell & $\wt$ in Ch.3 \S4, $Y_t$ in Ch.3 \S7 \\
$\tau$ & burst time & same \\
$\Kb$ & burst size (random) & $\mathcal K$ in Ch.3 \S3, $K$ in Ch.4 \S7--\S8 \\
$I(t)$ & $\Prob{\xt\neq R}$, no-burst probability & same \\
$D(t)$ & $\Prob{\text{cleared internally, no burst}}$ & \emph{not named}; equals $p_0(t)$ of Ch.4 \S5 \\
$\Ihat(t)$ & $\Prob{\xt\notin\{0,R\}}$, non-fixation & $\hat I$; \textbf{formula wrong for $\mu>0$} \\
$J(t)$ & $\E[\xt]$ & same \\
$K(t)$ & $\E[\xt^2]$ & same \\
$V(t)$ & $\E[\wt]$ & same \\
$\theta$ & $\beta(a-b)$, the characteristic exponent & written out longhand everywhere \\
$\Aa,\Bb$ & $a-1$ and $1-b$ & not used; they simplify almost everything \\
\midrule
$T$ & target-cell count (BMVR) & $S$ or $T$ \\
$\gamma$ & infection rate (BMVR) & $\beta$ --- \textbf{collides with birth rate} \\
$\Icell(t)$ & infected-cell count (BMVR) & $I$ --- \textbf{collides with $I(t)$} \\
$\Vfree(t)$ & free-virion count (BMVR) & $V$ --- \textbf{collides with $V(t)$} \\
$p$ & per-cell release rate (BMVR) & same \\
$c$ & virion clearance rate (BMVR) & same \\
$d_{\Icell}$ & infected-cell removal rate (BMVR) & $\delta$ --- \textbf{collides with catastrophe rate} \\
$r$ & population growth rate of the infection & not used \\
\bottomrule
\end{tabular}
\end{center}

%%%%%%%%%%%%%%%%%%%%%%%%%%%%%%%%%%%%%%%%%%%%%%%%%%%%%%%%%%%%%%%%%%%%%%%%%%%%%%%
\part*{Part I --- The single-cell BDC process, consolidated}
\addcontentsline{toc}{section}{\textbf{Part I --- The single-cell BDC process, consolidated}}
%%%%%%%%%%%%%%%%%%%%%%%%%%%%%%%%%%%%%%%%%%%%%%%%%%%%%%%%%%%%%%%%%%%%%%%%%%%%%%%

\section{Process and root algebra}

One cell, one BDC process, $X_0=1$ unless stated. Per-capita rates: birth $\beta$,
death $\mu$, catastrophe $\delta$. Catastrophe sends the process to the isolated
absorbing state $R$ and simultaneously sets $\wt=X_{\tau^-}$.

\subsection{The roots}

\begin{equation}
\eta=\frac{\beta+\mu+\delta}{2\beta},\qquad
a=\eta+\sqrt{\eta^2-\tfrac{\mu}{\beta}},\qquad
b=\eta-\sqrt{\eta^2-\tfrac{\mu}{\beta}} .
\end{equation}

Write
\begin{equation}
\Aa:=a-1>0,\qquad \Bb:=1-b>0,\qquad \theta:=\beta(a-b)=\sqrt{(\beta+\mu+\delta)^2-4\beta\mu}.
\end{equation}

\begin{framed}
\textbf{Identities used constantly (all elementary, none stated in the thesis):}
\begin{align}
ab&=\frac{\mu}{\beta}, &
a+b&=\frac{\beta+\mu+\delta}{\beta}, \label{eq:vieta}\\
\Aa\Bb=(a-1)(1-b)&=\frac{\delta}{\beta}, &
\Aa+\Bb&=a-b. \label{eq:AB}
\end{align}
\end{framed}

Equation~\eqref{eq:AB} is the one that does the work: it is why $b<1<a$ always (their
product is positive and $\beta(1-a)(1-b)=-\delta<0$), and it converts almost every messy
expression in Chapter~3 into a short one. \textbf{It should be added to the reference table
in Ch.3 \S3.}

\subsection{The three fixation functions}

There are \emph{three} absorbing-behaviour probabilities, not two. Chapter~3 defines
two and gets the second one wrong.

\begin{align}
I(t)&=\Prob{\text{no burst by }t}, \label{eq:Idef}\\
D(t)&=\Prob{\text{internal extinction by }t,\ \text{no burst}}, \label{eq:Ddef}\\
\Ihat(t)&=\Prob{\xt\notin\{0,R\}}=I(t)-D(t). \label{eq:Ihatdef}
\end{align}

$I$ and $D$ \emph{both} satisfy the same Riccati equation, because both events
factorise over independent subtrees:
\begin{equation}
\label{eq:riccati}
\ddt{f}=\beta(f-a)(f-b)=\mu-(\beta+\mu+\delta)f+\beta f^2,
\end{equation}
with $I(0)=1$ and $D(0)=0$. Writing $w:=\ee^{\theta t}\ge1$:
\begin{empheq}[box=\fbox]{align}
I(t)&=\frac{a\Bb+b\Aa w}{\Bb+\Aa w}, &
D(t)&=\frac{ab\,(w-1)}{aw-b}, \label{eq:ID}\\[0.3em]
\Ihat(t)&=\frac{(a-b)\Bb}{\Bb+\Aa w}+\frac{(a-b)\,b}{aw-b}
\;=\;\frac{(a-b)^2\,w}{(\Bb+\Aa w)(aw-b)}. &&\label{eq:Ihat}
\end{empheq}

Limits: $I(\infty)=b$, $D(\infty)=b$, $\Ihat(\infty)=0$, $\Ihat(0)=1$,
$\Ihat'(0)=-(\mu+\delta)$.

\begin{framed}
\textbf{Correction (major).} Chapter~3 \S5 and \S6.2 give
$\Ihat(t)=\hat\eta/\bigl(1-(1-\hat\eta)\ee^{\hat\eta\beta t}\bigr)$ with
$\hat\eta=(\beta+\mu+\delta)/\beta$, derived from the claim
$\Ihat'=-(\beta+\mu+\delta)\Ihat+\beta\Ihat^2$. That branching term is wrong.
On a birth event the two subtrees must satisfy \emph{no burst in either, and not both
extinct} --- which is $I^2-D^2$, not $\Ihat^2$. The correct ODE is non-autonomous,
$\Ihat'=-(\beta+\mu+\delta)\Ihat+\beta\Ihat(\Ihat+2D)$, and the clean route is
\eqref{eq:Ihatdef}.

The error is invisible when $\mu=0$ (then $D\equiv0$ and $\Ihat=I$), which is presumably
why it survived. At $(\beta,\mu,\delta)=(1,0.2,0.05)$, $t=1$: the correct value is
$\Ihat=0.8087$, the thesis formula gives $0.6675$.

\textbf{Independent confirmation.} Chapter~4 \S5's state probabilities $p_n(t)$ sum
(over $n\ge1$) to exactly the right-hand side of \eqref{eq:Ihat} --- verified
algebraically. So Ch.4 already implies the correct $\Ihat$ and contradicts Ch.3.
\end{framed}

\subsection{Multiplicity of infection}

For $X_0=k$ independent founders:
\begin{equation}
I_k(t)=I(t)^k,\qquad D_k(t)=D(t)^k,\qquad
\boxed{\ \Ihat_k(t)=I(t)^k-D(t)^k\ }
\end{equation}
Chapter~3 \S5 asserts $\Ihat_k=\Ihat^{\,k}$ and appends ``(Is this correct?)''. It is not;
the above is. The instinct was right.

\section{Moments}

\subsection{Simplified closed forms}

From $I'=-\delta J$ and \eqref{eq:riccati}:

\begin{empheq}[box=\fbox]{align}
J(t)&=\E[\xt]=-\frac{1}{\delta}\ddt{I}
=\frac{(a-b)^2\,w}{(\Bb+\Aa w)^2}
=-\frac{\beta}{\delta}(I-a)(I-b), \label{eq:J}\\[0.4em]
K(t)&=\E[\xt^2]=\Bigl[1+\tfrac{2\beta}{\delta}(1-I)\Bigr]J
=\frac{2\beta}{\delta}(\kappa-I)J,\qquad \kappa=1+\frac{\delta}{2\beta},\label{eq:K}\\[0.4em]
V(t)&=\E[\wt]=(1-I)\Bigl[1+\tfrac{\beta}{\delta}(1-I)\Bigr].\label{eq:V}
\end{empheq}

\begin{framed}
\textbf{Simplifications not in the thesis.}
\begin{itemize}[leftmargin=1.4em,itemsep=0.15em]
\item Ch.3 \S5 writes
$J(t)=-\lambda(1-a)(1-b)(a-b)^2\ee^{\beta(a-b)t}\big/\bigl[\delta(\cdots)^2\bigr]$.
Since $-\beta(1-a)(1-b)=\delta$ by \eqref{eq:AB}, the $\beta$ and $\delta$ cancel
identically and $J$ is just $(a-b)^2w/(\Bb+\Aa w)^2$, with $J(0)=1$ manifest.
\item Ch.3 gives $V=\frac{\beta-\mu}{\delta}(1-I)-J+1$. This equals
\eqref{eq:V}, which is a \emph{quadratic in $(1-I)$ alone} --- no $J$, no $\mu$, no
separate constant. Verified numerically at $(1,0.2,0.05)$, $t=1$: both give $0.17843$.
\item Consequently $\delta K=V'=\frac{\beta}{\delta}\frac{\dd}{\dd t}(I-\kappa)^2$, i.e.
$V=\frac{\beta}{\delta}\bigl[(I-\kappa)^2-(1-\kappa)^2\bigr]$, a useful alternative.
\end{itemize}
\end{framed}

\subsection{Late-time release and the mean burst size}

\begin{empheq}[box=\fbox]{equation}
V_\infty=\Bb\Bigl(1+\frac{\beta}{\delta}\Bb\Bigr)=\frac{a(1-b)}{a-1}=\frac{a\Bb}{\Aa}.
\label{eq:Vinf}
\end{empheq}

The middle form is Ch.3's $\frac{\beta-\mu}{\delta}(1-b)+1$ rewritten; the right-hand form
is new and much the most useful. Verified at $(1,0.2,0.05)$: both give $13.9857$.

$V_\infty$ is $\E[\Kb\mathbf 1\{\text{burst}\}]$ --- it \emph{includes} the realisations
that cleared internally and released nothing. The conditional mean burst size is
\begin{equation}
\E[\Kb\mid\text{burst}]=\frac{V_\infty}{1-b}=\frac{a}{a-1}
=\frac{\beta-\mu}{\delta}+\frac{1}{1-b}.
\end{equation}
The last equality is Ch.3 \S5's formula; the middle one is new and trivial by
\eqref{eq:Vinf}. Chapter~4 \S8 flags this distinction with ``(check)'' --- it is now
checked, and both forms agree.

\section{Burst time, burst size, and the quasi-stationary distribution}

\subsection{Burst time}

\begin{empheq}[box=\fbox]{equation}
\Prob{\tau>t}=I(t),\qquad
\varphi(t):=\text{density of }\tau=-I'(t)=\delta J(t),\qquad
\int_0^\infty\varphi=1-b.
\end{empheq}

This is defective: mass $b$ sits at $\tau=\infty$. Chapter~4 \S9 refers to
``the known pdf of rupture times $\phi(t)$'' without ever writing it down; it is
$\delta J(t)$, and it has been available since $I'=-\delta J$ in Chapter~3 \S7.

\subsection{Burst size is exactly the quasi-stationary distribution}

Chapter~4 \S5 gives, for $n\ge1$,
\begin{equation}
\label{eq:pn}
p_n(t)=\Bigl(\frac{a-b}{b-aw}\Bigr)^{\!2}w\;P(t)^{\,n-1},
\qquad P(t):=\frac{1-w}{b-aw},\qquad w=\ee^{\theta t}.
\end{equation}
The $n$-dependence is entirely in $P^{n-1}$, so \emph{at every time $t$}, conditional on
$\xt\ge1$, the load is geometric with ratio $P(t)$; and $P(t)\to 1/a$. Hence:

\begin{empheq}[box=\fbox]{equation}
\text{QSD}(n)=\lim_{t\to\infty}\Prob{\xt=n\mid \xt\notin\{0,R\}}
=\Bigl(1-\tfrac1a\Bigr)a^{-(n-1)}=(a-1)\,a^{-n},\quad n\ge1.
\end{empheq}

\begin{empheq}[box=\fbox]{equation}
\Prob{\Kb=k\mid\text{burst}}=\frac{(\delta/\beta)a^{-k}}{1-b}=(a-1)a^{-k},\quad k\ge1.
\end{empheq}

\begin{framed}
\textbf{New identity.} \emph{The burst-size distribution is exactly the quasi-stationary
distribution.} Both are geometric on $\{1,2,\dots\}$ with ratio $1/a$. This is not
obvious --- bursts are size-biased draws from the transient load distribution, integrated
against the burst intensity --- and it deserves a conceptual proof
(Task~\ref{task:qsdburst}). It is a striking, quotable exact result and neither chapter
notices it.
\end{framed}

Immediate consequences, none of which appear in the thesis:
\begin{align}
\langle X\rangle_{\rm QS}&=\frac{a}{a-1}=\frac{a}{\Aa}, &
\langle X^2\rangle_{\rm QS}&=\frac{a(a+1)}{(a-1)^2}, &
\mathrm{Var}_{\rm QS}(X)&=\frac{a}{(a-1)^2}. \label{eq:QSmoments}
\end{align}
Note $\langle X\rangle_{\rm QS}=\E[\Kb\mid\text{burst}]$ and, when $\mu=0$ only,
both coincide with $V_\infty$ --- which is precisely the coincidence that Chapter~4 \S6
observed and that Chapter~4 \S9 then over-generalised. For $\mu>0$ they differ:
at $(1,0.2,0.05)$, $\langle X\rangle_{\rm QS}=17.232$ while $V_\infty=13.986$.

\subsection{Size-biasing and the size of a late burst}

A cell bursts at rate $\delta\xt$, so the burst size given burst at time $t$ is a
\emph{size-biased} draw from $p_\cdot(t)$:
\begin{equation}
\Prob{\Kb=k\mid\tau=t}=\frac{k\,p_k(t)}{J(t)},\qquad
\E[\Kb\mid\tau=t]=\frac{K(t)}{J(t)}.
\end{equation}
This is exactly the quantity Chapter~4 \S9 calls $\E(R\mid\tau=t)$ and leaves undefined.
Its late-time limit is
\begin{equation}
\lim_{t\to\infty}\E[\Kb\mid\tau=t]=\frac{2\beta}{\delta}(\kappa-b)=1+\frac{2\beta}{\delta}(1-b),
\end{equation}
so a very late burst is roughly twice the size of an average one --- the signature of
size-biasing a geometric.

\subsection{Resolving the ``curious circularity'' of Chapter 4 \S8}

Chapter~4 \S8 computes $\E[\Kb\mid\tau>t]=(\E[\Kb]-V(t))/I(t)$, finds that its
$t\to\infty$ limit returns $\E[\Kb]$, and records that this is not understood.

The circularity is an artefact of using $\E[\Kb]$ to mean $\E[\Kb\mid\text{burst}]$
in the numerator while $I(t)$ in the denominator refers to the unconditioned event.
With consistent definitions:
\begin{equation}
\E\bigl[\Kb\mathbf 1\{\text{burst}\}\mid\tau>t\bigr]=\frac{V_\infty-V(t)}{I(t)}
\;\xrightarrow[t\to\infty]{}\;\frac{V_\infty-V_\infty}{b}=0,
\end{equation}
which is correct and unmysterious: conditioning on ``has not burst yet'' converges to
``never bursts'', on which the burst size is zero.

The genuinely interesting limit is obtained by conditioning on eventual burst,
\begin{equation}
\E[\Kb\mid \tau>t,\ \text{burst eventually}]=\frac{V_\infty-V(t)}{I(t)-b}
\;\xrightarrow[t\to\infty]{}\;\frac{K(t)}{J(t)}\Big|_{t\to\infty}
=\frac{\langle X^2\rangle_{\rm QS}}{\langle X\rangle_{\rm QS}},
\end{equation}
by l'H\^opital, since $V'=\delta K$ and $I'=-\delta J$. \textbf{This should replace
the closing paragraph of Ch.4 \S8.}

\subsection{The burst-size law is now proved, not just conjectured}

Chapter~4 \S7 states $\Prob{\Kb=k}=\frac{\delta}{\beta}a^{-k}$, attributes the
antiderivative to Wolfram Alpha, and says ``The current goal is to verify this using
Gillespie algorithm simulations''. No simulation is needed; the following is a proof.

With $P(t)=(1-w)/(b-aw)$ and $w=\ee^{\theta t}$, direct differentiation gives
\begin{equation}
\frac{\dd}{\dd t}\Bigl[\frac{P^k}{k\beta}\Bigr]
=\frac{1}{\beta}P^{k-1}P'
=\Bigl(\frac{a-b}{b-aw}\Bigr)^{\!2}w\,P^{\,k-1}=p_k(t),
\end{equation}
using $P'=\beta(a-b)^2w/(b-aw)^2$. Hence
\begin{equation}
\Prob{\Kb=k}=\delta k\!\int_0^\infty\! p_k
=\frac{\delta}{\beta}\bigl[P^k\bigr]_0^\infty
=\frac{\delta}{\beta}\Bigl(\frac1a\Bigr)^{\!k},
\end{equation}
since $P(0)=0$ and $P(\infty)=1/a$. \emph{Two} independent consistency checks now close
the loop:
\begin{enumerate}[leftmargin=1.6em,itemsep=0.15em]
\item Normalisation: $\sum_{k\ge1}\frac{\delta}{\beta}a^{-k}=\frac{\delta/\beta}{a-1}=1-b$
by \eqref{eq:AB}, matching the eventual-burst probability.
\item First moment: $\sum_{k\ge1}k\frac{\delta}{\beta}a^{-k}
=\frac{\delta}{\beta}\frac{a}{(a-1)^2}=\frac{a\Bb}{\Aa}=V_\infty$,
matching \eqref{eq:Vinf} --- verified algebraically for general $\mu$.
\end{enumerate}
The second check is the important one: it independently validates the burst law,
$V_\infty$, and the identity \eqref{eq:AB} simultaneously.

\section{Second moment of the release, corrected}

Chapter~3 \S8 derives $\E[\wt^2]$ and makes two errors that propagate into
$\mathrm{Var}(\wt)$ and into Ch.3 \S9.

Starting from $\dd\E[\wt^2]/\dd t=\delta\E[\xt^3]$ and
$K'=2(\beta-\mu)K+(\beta+\mu)J-\delta\E[\xt^3]$, eliminating $\E[\xt^3]$ and using
$J=-\delta^{-1}I'$ (note the sign) and $K=\delta^{-1}V'$:
\begin{equation}
\ddt{\E[\wt^2]}=\frac{2(\beta-\mu)}{\delta}\ddt{V}-\frac{\beta+\mu}{\delta}\ddt{I}-\ddt{K}.
\end{equation}
Integrating with the correct initial conditions
$\bigl(I,J,K,V,\E[W^2]\bigr)(0)=(1,1,\mathbf 1,0,0)$ --- note $K(0)=\E[X_0^2]=1$,
\emph{not} $0$ as the thesis states ---
\begin{empheq}[box=\fbox]{equation}
\E[\wt^2]=\frac{2(\beta-\mu)}{\delta}V-\frac{\beta+\mu}{\delta}I-K+\frac{\beta+\mu}{\delta}+1,
\label{eq:EW2}
\end{empheq}
\begin{equation}
\mathrm{Var}(\wt)=\E[\wt^2]-V^2 .
\end{equation}

\begin{framed}
\textbf{Errors corrected.} Chapter~3 \S8 eq.~(meanEY2) reads
$\E[Y_t^2]=\frac{2(\beta-\mu)}{\delta}V+\frac{\beta+\mu}{\delta}I+K-\frac{\beta+\mu}{\delta}$:
the signs on $I$ and $K$ are both wrong (traceable to the prose ``using
$J=\delta^{-1}\dd I/\dd t$'', which should be $J=-\delta^{-1}\dd I/\dd t$), and the
constant $+1$ is missing because $K(0)$ was taken as $0$.

\textbf{Verification of \eqref{eq:EW2}.} Take $\mu=0$, $t\to\infty$: $I,K\to0$,
$V\to1+\beta/\delta$, giving
$\E[W_\infty^2]=2\beta^2/\delta^2+3\beta/\delta+1$.
Independently, $W_\infty=\Kb\sim$ geometric on $\{1,2,\dots\}$ with success probability
$s=\delta/(\beta+\delta)$, whose second moment is $(2-s)/s^2=(2\beta+\delta)(\beta+\delta)/\delta^2
=2\beta^2/\delta^2+3\beta/\delta+1$. Exact agreement. The thesis formula gives a different
number.
\end{framed}

\section{Summary table of single-cell results}

\begin{center}
\renewcommand{\arraystretch}{1.45}
\begin{tabular}{@{}lll@{}}
\toprule
Quantity & Closed form & Status \\
\midrule
$I(t)$ & $(a\Bb+b\Aa w)/(\Bb+\Aa w)$ & thesis, correct \\
$D(t)$ & $ab(w-1)/(aw-b)$ & thesis (as $p_0$), not named \\
$\Ihat(t)$ & $(a-b)^2w/[(\Bb+\Aa w)(aw-b)]$ & \textbf{corrected here} \\
$\Ihat_k(t)$ & $I^k-D^k$ & \textbf{corrected here} \\
$J(t)$ & $(a-b)^2w/(\Bb+\Aa w)^2$ & simplified here \\
$K(t)$ & $\bigl[1+\tfrac{2\beta}{\delta}(1-I)\bigr]J$ & thesis, correct \\
$V(t)$ & $(1-I)\bigl[1+\tfrac{\beta}{\delta}(1-I)\bigr]$ & simplified here \\
$V_\infty$ & $a(1-b)/(a-1)$ & simplified here \\
$\E[\wt^2]$ & eq.~\eqref{eq:EW2} & \textbf{corrected here} \\
$\varphi(t)$ (burst time) & $\delta J(t)$, mass $1-b$ & implicit only \\
$\Prob{\Kb=k}$ & $(\delta/\beta)a^{-k}$ & thesis; \textbf{proved here} \\
QSD & geometric$(1/a)$ on $\{1,2,\dots\}$ & \textbf{new} \\
$\langle X\rangle_{\rm QS}$ & $a/(a-1)$ & \textbf{new} \\
$\E[\Kb\mid\tau=t]$ & $K(t)/J(t)$ & \textbf{new} \\
$\int_0^\infty\Ihat$ & $\beta^{-1}\log\bigl(a/(a-1)\bigr)$ & \textbf{new}, see \S\ref{sec:lifetime} \\
\bottomrule
\end{tabular}
\end{center}

%%%%%%%%%%%%%%%%%%%%%%%%%%%%%%%%%%%%%%%%%%%%%%%%%%%%%%%%%%%%%%%%%%%%%%%%%%%%%%%
\part*{Part II --- The burst-aware BMVR}
\addcontentsline{toc}{section}{\textbf{Part II --- The burst-aware BMVR}}
%%%%%%%%%%%%%%%%%%%%%%%%%%%%%%%%%%%%%%%%%%%%%%%%%%%%%%%%%%%%%%%%%%%%%%%%%%%%%%%

\section{Why the Chapter 4 \S9 attempt cannot be repaired as written}

Section 9 proposes replacing the constant $p$ in $\dd\Vfree/\dd t=p\Icell-c\Vfree$ by
$\langle X\rangle_{\rm QS}$.

\begin{framed}
\textbf{The obstruction is dimensional.} $p$ has units of particles per unit time.
$\langle X\rangle_{\rm QS}$ is a pure number. The proposed substitution therefore
introduces a spurious factor of time.

Worse, by \S\ref{sec:lifetime} below, when $\mu=0$ one has
$\langle X\rangle_{\rm QS}=V_\infty$ exactly --- so the proposal substitutes the
\emph{total lifetime output of one cell} in place of its \emph{output rate}.
This is why the section stalls, and no reweighting of a count by a probability density
can fix it.
\end{framed}

The two candidate integrals in \S9 fail for related reasons:
\begin{itemize}[leftmargin=1.6em,itemsep=0.2em]
\item $p=\int\varphi(t)\,\E[\Kb\mid\tau=t]\,\dd t$ evaluates, using
$\varphi=\delta J$ and $\E[\Kb\mid\tau=t]=K/J$, to $\delta\int K=V_\infty$ --- again a
count, not a rate.
\item $p=\delta\int\varphi(t)\,\E[\xt^2\mid\cdot]\,\dd t$ has the right dimensions
\emph{if} $\int\varphi\,\dd t$ is read as an averaging operator, but weights by the
\emph{burst-time} density. The correct weight is the \emph{age distribution of currently
infected cells}, which is a different measure entirely.
\end{itemize}

Section 9 also asserts $\Vfree(t)=N_0V(t)$ for $N_0$ initially infected cells. That is
correct only if all $N_0$ cells are infected at the same instant. The section then
notices the problem (``at any time, there will be infected cells that begin their
infection dynamics at a number of different start times'') but does not act on it. Acting
on it is the whole solution.

\section{The correct object: a renewal (age-structured) BMVR}

\subsection{Kernels}

Two kernels, both already in closed form from Part I:
\begin{empheq}[box=\fbox]{align}
\text{cell-survival kernel:}\quad &\Ihat(a)=\Prob{\text{cell infected }a\text{ ago is
still productively infected}},\\
\text{release kernel:}\quad &g(a)=\delta K(a)=V'(a)=\text{expected release flux at
cell-age }a.
\end{empheq}
$g$ is the correct object because a cell bursts at rate $\delta\xt$ releasing $\xt$
particles, so the expected release in $[a,a+\dd a]$ is $\delta\E[\xt^2]\dd a$.
Chapter~4 \S9 states this fact in its second paragraph and then does not use it.

\subsection{The system}

With $T$ target cells (held constant, as in the simplified BMVR), infection rate
$\gamma$, virion clearance $c$, and incidence $i(t)=\gamma T\Vfree(t)$:
\begin{empheq}[box=\fbox]{align}
i(t)&=\gamma T\,\Vfree(t),\\
\Icell(t)&=\int_0^t i(t-\alpha)\,\Ihat(\alpha)\,\dd\alpha,\label{eq:Icell}\\
\ddt{\Vfree}&=\int_0^t i(t-\alpha)\,\delta K(\alpha)\,\dd\alpha-c\,\Vfree(t).
\label{eq:Vfree}
\end{empheq}
(Initial-cohort terms are omitted for clarity; add
$\Icell_0\Ihat(t)$ and $\Icell_0\delta K(t)$ for cells present at $t=0$.)

This replaces \emph{both} BMVR terms: the constant release $p\Icell$ becomes a
convolution, and the constant cell removal $d_\Icell\Icell$ becomes the age-structured
form implicit in \eqref{eq:Icell}. Chapter~4 \S9 only ever proposed to change the first.

\subsection{Relation to the literature (to be stated honestly in any write-up)}

The renewal/age-structured formulation itself is standard (McKendrick--von Foerster;
Nelson--Gilchrist--Perelson-type age-structured within-host models; Euler--Lotka
reduction from demography). \textbf{The novelty claim is not the framework.} It is that
$\Ihat(a)$ and $g(a)$ are here \emph{derived in closed form from a mechanistic
intracellular stochastic process} rather than posited phenomenologically, together with
the consequences in \S\ref{sec:peff}--\S\ref{sec:flooding}. Any submission must say this
plainly in the introduction; claiming the framework will draw fire.

\section{Effective BMVR parameters}
\label{sec:peff}

\subsection{Mean productive lifetime}
\label{sec:lifetime}

\begin{empheq}[box=\fbox]{equation}
\E[T_{\rm prod}]=\int_0^\infty\Ihat(t)\,\dd t=\frac{1}{\beta}\log\!\Bigl(\frac{a}{a-1}\Bigr)
=\frac{1}{\beta}\log\langle X\rangle_{\rm QS}.
\end{empheq}

\emph{Proof.} Write $\Ihat=(I-b)-(D-b)$ with
$I-b=\Bb(a-b)/(\Bb+\Aa w)$ and $D-b=-b(a-b)/(aw-b)$; substitute $w=\ee^{\theta t}$,
use $\theta=\beta(a-b)$, and integrate by partial fractions in $w$:
\begin{equation}
\int_0^\infty\Ihat\,\dd t
=\frac{1}{\beta}\Bigl[\log\frac{a-b}{a-1}+\log\frac{a}{a-b}\Bigr]
=\frac{1}{\beta}\log\frac{a}{a-1}. \qquad\square
\end{equation}
When $\mu=0$ this is $\beta^{-1}\log(1+\beta/\delta)$. Checked numerically at
$(1,0.2,0.05)$: formula $2.8468$, trapezoidal integration $2.8541$ at $h=1$ (converging
from above as expected). The relation to $\langle X\rangle_{\rm QS}$ is a pleasant
one and worth remarking on: the productive lifetime is the time it takes exponential
growth at rate $\beta$ to reach the quasi-stationary level.

\subsection{Exact reduction to a constant-parameter BMVR}

Substitute $i(t)=i_0\ee^{rt}$ into \eqref{eq:Icell}--\eqref{eq:Vfree} and compare
with the BMVR equations
$\Icell'=i-d_\Icell\Icell$, $\Vfree'=p\Icell-c\Vfree$. Writing
$\Lap{f}(r)=\int_0^\infty\ee^{-r\alpha}f(\alpha)\dd\alpha$:

\begin{empheq}[box=\fbox]{align}
p_{\rm eff}(r)&=\frac{\delta\,\Lap{K}(r)}{\Lap{\Ihat}(r)},
&
d_{\Icell,{\rm eff}}(r)&=\frac{1}{\Lap{\Ihat}(r)}-r.
\label{eq:peff}
\end{empheq}

\begin{framed}
\textbf{This reduction is exact, not approximate, in the exponential phase.}
Substituting \eqref{eq:peff} into the BMVR eigenvalue relation
$(r+d_\Icell)(r+c)=\gamma Tp$ returns
$r+c=\gamma T\delta\Lap{K}(r)$, which is precisely the characteristic equation of the
renewal system. So for any $r$ there is a unique constant-parameter BMVR reproducing the
renewal model's exponential solution exactly. It is \emph{not} exact in transients.
\end{framed}

Limits, both checked:
\begin{align}
r=0:\quad & p_{\rm eff}=\frac{V_\infty}{\int_0^\infty\Ihat}
=\frac{\beta\,V_\infty}{\log\bigl(a/(a-1)\bigr)},
\qquad d_{\Icell,{\rm eff}}=\frac{\beta}{\log\bigl(a/(a-1)\bigr)},\\
r\to\infty:\quad & p_{\rm eff}\to\delta,\qquad d_{\Icell,{\rm eff}}\to\mu+\delta.
\end{align}
The $r=0$ case uses $\int_0^\infty\delta K=V_\infty$ (immediate from $V'=\delta K$),
which is the same answer as naive moment-matching --- a useful consistency check.
The $r\to\infty$ case is the correct young-cell limit: a newly infected cell holds one
particle, bursts at rate $\delta$, and releases one.

For $\mu=0$ everything is elementary:
\begin{equation}
p_{\rm eff}(0)=\frac{\beta(1+\beta/\delta)}{\log(1+\beta/\delta)}.
\end{equation}

\subsection{Closed forms for the transforms}

With $s:=r/\theta$ and $z:=-\Bb/\Aa$:
\begin{empheq}[box=\fbox]{align}
\Lap{\Ihat}(r)&=\frac{1}{\beta(s+1)}
\left[\frac{\Bb}{\Aa}\,\Fhyp\Bigl(1,s+1;s+2;-\frac{\Bb}{\Aa}\Bigr)
+\frac{b}{a}\,\Fhyp\Bigl(1,s+1;s+2;\frac{b}{a}\Bigr)\right],\\[0.4em]
\delta\Lap{K}(r)&=\frac{\delta}{\beta}
\left[\frac{\Fhyp(1,s;s+2;z)}{\Aa(s+1)}
+\frac{2\,\Fhyp(2,s;s+3;z)}{\Aa^2(s+1)(s+2)}\right].
\end{empheq}

Both follow from the substitution $v=\ee^{-\theta t}$, which maps the integrals onto
$\int_0^1v^{s-1}(1-v)^n(\Aa+\Bb v)^{-m}\dd v
=\Aa^{-m}B(s,n+1)\,\Fhyp(m,s;s+n+1;z)$. Since $z<0$ and $b/a<1$, both series converge
on the physical range. At $s=0$ both reduce, via $\Fhyp(1,1;2;z)=-\log(1-z)/z$, to
$\int\Ihat=\beta^{-1}\log(a/(a-1))$ and $\delta\Lap K(0)=V_\infty$ --- both checks done.

\begin{remark}
The appearance of $\Fhyp$ here and in Chapter~5's two-type reduction is not a
coincidence worth over-selling, but it does make the two chapters look like one project
rather than two.
\end{remark}

\section{Characteristic equation, $R_0$, and the growth-phase dependence of $p$}

\subsection{$R_0$ is untouched by bursting}

The characteristic equation of \eqref{eq:Icell}--\eqref{eq:Vfree} is
\begin{empheq}[box=\fbox]{equation}
r+c=\gamma T\,\delta\Lap{K}(r),\qquad\text{hence}\qquad
R_0=\frac{\gamma T\,V_\infty}{c}.
\end{empheq}

\begin{framed}
\textbf{Result of interest.} At fixed total per-cell output $V_\infty$, replacing
constant budding by bursting leaves $R_0$ \emph{exactly unchanged}. Bursting alters the
generation-time distribution --- hence $r$, hence the extinction probability, hence the
stochastic early phase --- but not the threshold. This is the classic $R_0$-versus-$r$
distinction, and it connects the BMVR extension directly to the ``perilous beginning''
theme of Ch.3 \S2 and to the flooding-advantage question of Ch.1 \S6.
\end{framed}

\subsection{$p$ is not a constant of the pathogen}

Since $p_{\rm eff}$ depends on $r$ and decreases from $p_{\rm eff}(0)$ towards $\delta$,
the effective release rate is smaller during acute exponential growth (young cells
dominate, loads are low) than at a steady state. At $(\beta,\mu,\delta)=(1,0,0.1)$:
$p_{\rm eff}(0)=4.59$ against $p_{\rm eff}(\infty)=0.1$.

\begin{framed}
\textbf{Prediction.} BMVR fits to acute-phase data and to set-point data should yield
systematically different $p$, in a direction and by an amount the theory specifies.
Equivalently: the eclipse-like delay that is normally inserted by hand into within-host
models emerges here from the intracellular dynamics, with no extra parameter.
\end{framed}

\subsection{Identifiability (a negative result worth publishing)}

Three intracellular rates $(\beta,\mu,\delta)$ map onto two BMVR parameters
$(p,d_\Icell)$ at any given $r$. The map is therefore not invertible.

\begin{framed}
\textbf{Consequence for Ch.4 \S9's stated aim.} The section's purpose is: ``given someone
\dots{} has a good estimate for $p$ \dots{} they may then be able to access an
approximation for the rates constituting the QS mean.'' They cannot --- not from $p$
alone, and not from $(p,d_\Icell)$ either. A third observable is required. The natural
candidates, in order of experimental accessibility:
\begin{enumerate}[leftmargin=1.6em,itemsep=0.1em]
\item the burst-size dispersion (the geometric ratio $1/a$ fixes one combination directly);
\item the fraction of infected cells that clear without bursting, which is $b$;
\item the shape --- not just the mean --- of the release-time distribution $\varphi=\delta J$.
\end{enumerate}
Any one of these plus $(p,d_\Icell)$ over-determines $(\beta,\mu,\delta)$, giving a
consistency test. This should be stated as a result, not buried.
\end{framed}

\section{The flooding advantage, answered exactly}
\label{sec:flooding}

This is the strongest single result available and it addresses Ch.1 \S6 and Ch.7's
Komarova discussion directly.

\subsection{Offspring distribution of one infected cell}

Assume each released particle independently either infects ($\gamma T$) or is cleared
($c$), so infects with probability $q=\gamma T/(\gamma T+c)$. With $y=1-q+qz$:
\begin{empheq}[box=\fbox]{equation}
G_{\rm off}(z)=b+\frac{\delta}{\beta}\,\frac{y}{a-y},
\end{empheq}
which satisfies $G_{\rm off}(1)=b+(1-b)=1$. Moments:
\begin{align}
m&=qV_\infty\ (=R_0\text{ when }\gamma T\ll c), &
\mathrm{Var}&=\frac{2q^2V_\infty}{a-1}+qV_\infty-q^2V_\infty^2 .
\end{align}

\subsection{Extinction probability in closed form}

$G_{\rm off}(z)=z$ is a \emph{quadratic} in $y$ with $y=1$ always a root. The other root
gives
\begin{empheq}[box=\fbox]{equation}
z_{\rm ext}^{\rm burst}=\frac{a(1-q+qb)-1+q}{q}=\frac{a-1}{q}+1-a(1-b).
\end{empheq}

\subsection{Comparison with constant budding at matched $R_0$}

Comparator: a BMVR cell with $\mathrm{Exp}(d_\Icell)$ lifetime producing at constant rate
$p$, matched so that $p/d_\Icell=V_\infty$ and with the same $q$. Its offspring number is
geometric with mean $R_0$, so $z_{\rm ext}^{\rm bud}=1/R_0$.

Writing $L:=a(1-b)=a-\mu/\beta$, so that $V_\infty=L/(a-1)$:

\begin{empheq}[box=\fbox]{equation}
z_{\rm ext}^{\rm burst}-z_{\rm ext}^{\rm bud}
=\bigl(L-1\bigr)\Bigl(\frac{1}{R_0}-1\Bigr)
=\bigl(a-\tfrac{\mu}{\beta}-1\bigr)\frac{1-R_0}{R_0}.
\label{eq:flood}
\end{empheq}

\begin{framed}
\textbf{Main result.} For $R_0>1$, bursting strictly lowers the extinction probability
relative to matched-$R_0$ budding if and only if $L>1$, i.e.
\begin{equation}
a>1+\frac{\mu}{\beta}
\quad\Longleftrightarrow\quad
\delta(\beta+\mu)>\beta\mu
\quad\Longleftrightarrow\quad
\boxed{\ \frac{1}{\delta}<\frac{1}{\beta}+\frac{1}{\mu}\ }
\end{equation}
When $L=1$ the two extinction probabilities are \emph{identical for every $q$}.
For $R_0\le1$ both models go extinct almost surely and the comparison is empty ---
note that \eqref{eq:flood} then changes sign, which is an artefact of both formulas
running outside $[0,1)$, not a reversal of the biology.

Interpretation: the flooding advantage is real, but it is bought against a risk. Bursting
concentrates output into one batch, which lowers offspring variance relative to the
geometric output of a budding cell; but it also carries the all-or-nothing hazard that
the intracellular population dies out first, with probability $b$, and the cell releases
nothing. The condition says the advantage survives exactly when the burst timescale
$1/\delta$ is shorter than the harmonic-type combination of the birth and death
timescales. When $\mu=0$ the condition holds for all $\delta>0$: bursting always wins.
\end{framed}

\emph{Numerical verification.} $(\beta,\mu,\delta)=(1,0,0.1)$, $q=0.2$:
$z^{\rm burst}=0.400$, $z^{\rm bud}=0.4545$, and \eqref{eq:flood} predicts
$-0.05455$ --- exact. $(1,0.9,0.1)$, $q=0.9$: $z^{\rm burst}=0.9351$,
$z^{\rm bud}=0.8442$ --- the ordering reverses, as the condition requires.
$(1,0.5,1/3)$ sits exactly at $L=1$ and gives $z^{\rm burst}=z^{\rm bud}$ at both
$q=0.6$ and $q=0.9$.

\begin{remark}
This deserves care in the write-up: the comparison holds the offspring \emph{mean} fixed
and is therefore a statement about branching structure only. It says nothing about
nonlinear effects --- immune saturation, local depletion --- which is where Komarova's
own ``all run out at once'' intuition really lives. Both mechanisms should be named and
separated.
\end{remark}

\section{Multiplicity of infection}

For a cell founded by $k$ particles, the kernels are
\begin{align}
\Ihat_k(a)&=I(a)^k-D(a)^k, &
g_k(a)&=\delta\bigl[k\,K(a)+k(k-1)J(a)^2\bigr],
\end{align}
the second from the variance of a sum of $k$ i.i.d.\ copies. The release kernel is
\emph{superlinear} in $k$, so MOI matters even in the mean --- a genuine prediction that
the BMVR cannot express, and one that single-cell or low-MOI assays can test.

\section{Two-type extension and the eclipse phase}

Chapter~5 already provides the exact finite-time no-catastrophe probability $S(t)$ for
the two-type process with type-specific rates $(\delta_1,\delta_2)$.

\begin{framed}
\textbf{Observation not made in Chapter 5: GATE \emph{is} an eclipse phase.}
The regime $\delta_1=0<\delta_2$ means no burst is possible before conversion, so
$S'(0)=0$ --- a zero initial hazard, followed by a rising one. That is precisely the
eclipse-phase signature that within-host models normally insert by hand. In the two-type
BDC it emerges mechanistically, with the eclipse duration set by $\nu$.

This is the strongest reason to build the renewal model on the two-type kernel rather
than the one-type kernel, and it makes Chapters 4 and 5 one story instead of two.
\end{framed}

What is needed, and is \emph{not} yet available:
\begin{itemize}[leftmargin=1.6em,itemsep=0.2em]
\item The release kernel is
$g(a)=\delta_1\E[X_a^2]+(\delta_1+\delta_2)\E[X_aY_a]+\delta_2\E[Y_a^2]$
(all expectations over the killed process, i.e.\ including the indicator
$\mathbf 1\{\tau_c>a\}$). These second moments are not in Chapter~5.
\item The cell-survival kernel must be the two-type analogue of $\Ihat$, i.e.\ ``neither
burst nor extinct''. Chapter~5 gives only $S$, the analogue of $I$, and explicitly notes
that ordinary death does not stop its clock. The analogue of $D$ is missing.
\end{itemize}
See Tasks \ref{task:twotypemoments} and \ref{task:twotypeD}.

\section{Confrontation with single-cell data}

The Hataye-type assay (one infected cell per well, viral counts hourly) measures the
single-cell BDC directly, with no population model interposed. Predictions, all
one-parameter or parameter-free given $(\beta,\mu,\delta)$:

\begin{enumerate}[leftmargin=1.6em,itemsep=0.2em]
\item A fraction $b$ of wells should show \emph{no release ever}.
\item Cumulative release per well is a step function: zero, then a single jump.
\item Jump times follow density $\delta J(t)/(1-b)$.
\item Jump sizes follow geometric$(1/a)$ --- and by the identity of \S4.2, this is the
same law as the quasi-stationary load.
\item Between-well variance is $\mathrm{Var}(W_\infty)$ from \eqref{eq:EW2}, which the
deterministic BMVR cannot produce at all.
\end{enumerate}

\begin{framed}
\textbf{Honest caveat that must be in the paper.} HIV buds; it does not lyse in the
all-or-nothing way the BDC assumes. Predictions (2) and (3) are therefore too strong for
HIV as literally stated, and observed ``intermittency'' is not the same as a single
terminal burst. Either (i) frame the one-type BDC as the bursting endpoint of a spectrum
and fit an intermediate release model, which is exactly the budding--bursting spectrum of
Chapter~6, or (ii) apply it to a genuinely lytic system (\textit{Y.~pestis},
\textit{F.~tularensis}, lytic phage) where the assumption is right, and use HIV only as
motivation. Option (ii) is safer; option (i) is the better paper if Chapter~6 can supply
the intermediate model.
\end{framed}

%%%%%%%%%%%%%%%%%%%%%%%%%%%%%%%%%%%%%%%%%%%%%%%%%%%%%%%%%%%%%%%%%%%%%%%%%%%%%%%
\part*{Part III --- Open tasks}
\addcontentsline{toc}{section}{\textbf{Part III --- Open tasks}}
%%%%%%%%%%%%%%%%%%%%%%%%%%%%%%%%%%%%%%%%%%%%%%%%%%%%%%%%%%%%%%%%%%%%%%%%%%%%%%%

\section{Ordered by priority}

\begin{task}[Literature check --- do this first]
\label{task:lit}
Establish whether Carruthers et al.\ (\textit{F.~tularensis}) already couple an
intracellular BDC to a between-cell model. If they do, the novelty claim must shift
entirely onto the closed-form kernels, $p_{\rm eff}(r)$, and \eqref{eq:flood}. Also
check: Pearson--Krapivsky--Perelson on early stochastic infection with burst versus
budding; Nelson--Gilchrist--Perelson age-structured within-host; Gilchrist--Coombs
nested models; the eclipse-phase fitting literature; Heffernan--Wahl on burst size.
Confirm the BMVR attribution --- Ch.1 cites McLean 1993, but Nowak--Bangham 1996 and
Perelson et al.\ 1996 are the usual sources.
\end{task}

\begin{task}[Prove the QSD/burst-size identity]
\label{task:qsdburst}
Find a conceptual proof that $\Prob{\Kb=k\mid\text{burst}}$ equals the quasi-stationary
distribution. The algebra is done; a structural argument (size-biasing composed with the
burst intensity, or a time-reversal/Yaglom argument) would make it quotable rather than
coincidental.
\end{task}

\begin{task}[Two-type killed second moments]
\label{task:twotypemoments}
Derive $\E[X_a^2\mathbf 1_{\tau_c>a}]$, $\E[X_aY_a\mathbf 1]$, $\E[Y_a^2\mathbf 1]$ for
the Chapter~5 process. The moment hierarchy does not close upwards, exactly as in the
one-type case; the one-type escape was to obtain $K$ by differentiating the Riccati
equation. The two-type analogue should be to differentiate the backward system
\eqref{eq:S-two}--(G) below. This is the main technical obstacle to the two-type
renewal model and should not be assumed easy.
\end{task}

\begin{task}[Two-type non-fixation probability]
\label{task:twotypeD}
Construct the analogue of $D(t)$ for the two-type process (extinction without
catastrophe) and hence the two-type $\Ihat$. Chapter~5 supplies only $S$ and $G$.
\end{task}

\begin{task}[Growth-phase dependence, quantified]
Plot $p_{\rm eff}(r)$ and $d_{\Icell,{\rm eff}}(r)$ over the physiological range of $r$
for representative $(\beta,\mu,\delta)$, and quantify the predicted discrepancy between
acute-phase and set-point BMVR fits. This is the figure that will sell the paper.
\end{task}

\begin{task}[Characterise the flooding boundary]
Map the region $\delta(\beta+\mu)>\beta\mu$ against plausible parameter ranges for
\textit{Y.~pestis} in macrophages, and determine whether real pathogens sit on the
advantageous side. Combine with \eqref{eq:flood} to give the \emph{size} of the advantage.
\end{task}

\begin{task}[Population-level variance]
Propagate $\mathrm{Var}(\wt)$ through the branching structure to a population-level
prediction. For the single-cell assay this is unnecessary (\S9 of Part II), but for
whole-host data it is the main stochastic signal the BMVR cannot reproduce.
\end{task}

\begin{task}[Intermediate release models]
If option (i) of the data section is taken, define a release model interpolating between
pure budding and pure bursting --- e.g.\ partial release of a fraction of the load at
catastrophe, or a per-capita export rate alongside the catastrophe rate --- and check
whether the kernel remains closed-form. This is the natural link to Chapter~6.
\end{task}

\begin{task}[Nonlinear intracellular growth]
A referee will object that linear birth lets the conditional load grow without bound.
Prepare the defence (the burst hazard is proportional to load and so regulates it; the
QSD result is precisely the statement that the conditional load stabilises at
$a/(a-1)$), and check numerically how far a logistic intracellular birth rate moves
$g(a)$ and $\Ihat(a)$.
\end{task}

For reference, the Chapter~5 backward system to be differentiated in
Task~\ref{task:twotypemoments}:
\begin{align}
\ddt{S}&=\lambda_1S^2-(\lambda_1+\mu_1+\nu+\delta_1)S+\mu_1+\nu G,\label{eq:S-two}\\
\ddt{G}&=\lambda_2G^2-(\lambda_2+\mu_2+\delta_2)G+\mu_2,\qquad S(0)=G(0)=1.
\end{align}

%%%%%%%%%%%%%%%%%%%%%%%%%%%%%%%%%%%%%%%%%%%%%%%%%%%%%%%%%%%%%%%%%%%%%%%%%%%%%%%
\part*{Part IV --- Errata and incomplete items}
\addcontentsline{toc}{section}{\textbf{Part IV --- Errata and incomplete items}}
%%%%%%%%%%%%%%%%%%%%%%%%%%%%%%%%%%%%%%%%%%%%%%%%%%%%%%%%%%%%%%%%%%%%%%%%%%%%%%%

Tags: \ERR\ mathematical error \quad \GAP\ unfinished \quad \NOT\ notation collision
\quad \TYP\ typographical \quad \STR\ structural.

\section{Chapter 1 --- Introduction}
\emph{Folder:} \texttt{1 Intro/document MAIN/}

\begin{enumerate}[leftmargin=2.2em,label=\textbf{1.\arabic*}]
\item \ERR\ \texttt{sections/06\_hiv\_and\_the\_basic\_model\_of\_viral\_replication.tex}.
The BMVR is written as
$\dot S=\lambda-\beta SI$, $\dot I=\beta SI-\delta I$, $\dot V=pI-cV$.
The infection term should be driven by $V$, not $I$: $\dot I=\beta SV-\delta I$.
As written, $V$ is entirely decoupled and plays no role in the dynamics --- which
undermines the whole section, since the point is that $V$ is what the model gets wrong.
Chapter~4 \S9 uses the correct $\beta TV$ form, so the two are inconsistent with each
other.

\item \ERR\ Same file, simplified system: $\dot I=\beta TI-\delta I$, $\dot V=pI-cV$.
Same problem; should be $\dot I=\beta TV-\delta I$.

\item \NOT\ Same file. $\lambda$ is the target-cell influx rate here but the
intracellular \emph{birth} rate in Ch.3. $\delta$ is the infected-cell death rate here
but the \emph{catastrophe} rate in Ch.3/4/5. $\beta$ is the infection rate here but the
birth rate in Ch.3. All three collide at exactly the point where the two models are
supposed to meet. This is the single most damaging notational problem in the thesis.

\item \NOT\ $I$ and $V$ denote cell and virion \emph{counts} here, and \emph{probabilities
and expectations} in Ch.3/4. Ch.4 \S9 writes both meanings within a page.

\item \GAP\ Attribution: the BMVR is cited to \texttt{mclean1993balance}. Verify against
Nowak \& Bangham (1996) and Perelson et al.\ (1996), which are the conventional sources.

\item \GAP\ Verify the Hataye et al.\ experimental details as stated (``2018'',
``Los Alamos'', ``32 small wells'', hourly sampling) against the published paper; the
citation key is \texttt{hataye2019principles}, so the year in the text and the key
disagree.

\item \GAP\ \texttt{sections/08\_research\_questions.tex} lines 18--19. Research question
3 promises ``a potential update to the standard BMVR''. As things stand the promise is
not delivered --- Ch.4 \S9 states its own formula is incorrect. Either deliver
(Part II above) or soften the claim.

\item \NOT\ \texttt{main.tex} line 85 refers to ``the basic model of viral replication
(BMVR)'' as a defined acronym before \S6 introduces it. Check ordering.
\end{enumerate}

\section{Chapter 3 --- BDC core}
\emph{Folder:} \texttt{3 BDC core/document MAIN/}

\begin{enumerate}[leftmargin=2.2em,label=\textbf{3.\arabic*}]
\item \ERR\ \textbf{(major)} \fp{sections/05_analytical_quantities.tex} and
\fp{sections/06_derivation_of_equations_and_results.tex}, \S6.2.
The formula $\Ihat(t)=\hat\eta/\bigl(1-(1-\hat\eta)\ee^{\hat\eta\beta t}\bigr)$ and the
ODE $\Ihat'=-(\beta+\mu+\delta)\Ihat+\beta\Ihat^2$ are both wrong for $\mu>0$.
See Part I \S3.2 for the correct $\Ihat=I-D$ and for numerical evidence. Correct only when
$\mu=0$.

\item \ERR\ \texttt{sections/05\_analytical\_quantities.tex}.
$\Ihat_k(t)=[\Ihat(t)]^k$ is wrong; correct is $I^k-D^k$. The text itself appends
``(Is this correct?)''.

\item \ERR\ \texttt{sections/08\_variance\_and\_standard\_deviation\_of.tex},
eq.~(meanEY2). Signs on $I$ and $K$ are wrong and the constant $+1$ is missing.
Two causes: the prose states $J=\delta^{-1}\dd I/\dd t$ where the sign should be negative;
and the listed initial conditions give $K(0)=0$ when in fact $K(0)=\E[X_0^2]=1$.
Corrected formula and an independent geometric-moment check in Part I \S5.

\item \ERR\ Same file, eq.~(VarY): inherits the above.

\item \ERR\ \texttt{sections/09\_variable\_dependencies\_on.tex}.
$K(I)=-\frac{2\beta^2}{\delta}(I-\kappa)(I-I_-)(I-I_+)$ contradicts
\texttt{11\_variable\_dependencies\_on.tex}, which gives
$K(I)=\frac{2\beta^2}{\delta^2}(I-k)(I-I_+)(I-I_-)$. \textbf{The \S11 version is correct};
\S9 has both a wrong power of $\delta$ and a sign error. Derivation:
$K=\frac{2\beta}{\delta}(\kappa-I)J$ with $J=-\frac{\beta}{\delta}(I-a)(I-b)$.

\item \ERR\ \texttt{sections/06\_derivation\_of\_equations\_and\_results.tex}, eq.~(ddIIfact):
$\ddt{I}=\beta(I-I_+)(I-I_+)$ --- the subscript is repeated; should be $(I-I_+)(I-I_-)$.

\item \ERR\ \texttt{sections/02\_chapter\_introduction.tex}: ``The importance of
stochasticity can not be understated'' inverts the intended meaning. Should be
``overstated'' (or rephrase).

\item \GAP\ \texttt{sections/05\_analytical\_quantities.tex}: the placeholder
``PLOT $\E[\mathcal K]$ against $\delta$ for $\mu=0$, $\lambda=1$'' is still in the source.

\item \GAP\ \texttt{sections/05\_analytical\_quantities.tex}: $\E[\mathcal K]$ is
never cross-checked against the geometric burst law of Ch.4 \S7. They do agree; the check
is in Part I \S4.4 and should be added, since it validates several results at once.

\item \STR\ \textbf{Duplicated sections.} \texttt{08\_variance\_and\_standard\_deviation\_of.tex}
$\approx$ \texttt{10\_variance\_and\_standard\_deviation\_of.tex}, and
\texttt{09\_variable\_dependencies\_on.tex} $\approx$
\texttt{11\_variable\_dependencies\_on.tex}. All four are \verb|\input| in
\texttt{main.tex}, so the compiled PDF repeats the variance and $I$-dependency material
with two conflicting versions of $K(I)$. Sections 10--11 are the earlier, shorter drafts.
Decide which survives and delete the other pair.

\item \STR\ \texttt{sections/07\_expected\_number\_of\_units\_j\_t.tex} re-derives $J(t)$
and $V(t)$, already covered in \S5. Third statement of the same material.

\item \NOT\ Birth rate is $\lambda$ in the \S3 reference table and in the $I(t)$, $J(t)$
formulas of \S5, but $\beta$ from \S4 onwards. Two symbols, one rate, within one chapter.

\item \NOT\ The released count is $\wt$ in \S4 and the reference table, but $Y_t$ in \S7
and \S8.

\item \NOT\ $K$ is the second moment $\E[\xt^2]$ throughout, but $\mathcal K$ is the burst
size in the \S3 reference table, and Ch.4 \S7--\S8 use plain $K$ for the burst size.
Three uses of one letter.

\item \NOT\ The rupture state is $H$ here and $R$ in Ch.4 and Ch.5.

\item \NOT\ \S5 defines $\Ihat(t)=\Prob{\xt\notin\{0,H\}}$; \S6.2 defines
$\Ihat(t)=\Prob{\xt\neq0}$. These are different events. Fix to the former.

\item \TYP\ \texttt{sections/03\_reference\_values.tex}: ``$p_k(t)$: $\Prob{\xt=i}$'' ---
index mismatch, and likewise for $P_k$ and $\phi_k$. Also $\phi_k$ is called a
distribution with $\Prob{\tau=i}$, but $\tau$ is continuous: it is a density.

\item \TYP\ \texttt{sections/04\_process\_definition.tex}: the transition list omits the
conditioning on $\xt$ and the $o(\Delta t)$ terms present in \S6's careful version.
Also both figures carry \verb|\label{BDRsims}| --- duplicate label.

\item \TYP\ \texttt{sections/02\_chapter\_introduction.tex}: ``Brokwell'' and
``Brockwell'' both used; ``Jonathan Carrthurs'' $\to$ Carruthers; ``continuois'',
``macrophaegs'', ``mammels'', ``populance'', ``accross'', ``sppured'', ``Ricatti'' $\to$
Riccati, ``arives'', ``encorporated''.

\item \STR\ \texttt{sections/02\_chapter\_introduction.tex}: cross-references are stale.
``This will be explored in chapter 3'' and ``explored further in chapter 4'' point to
material now in Chapters 6 and 7 of the current tree.
\end{enumerate}

\section{Chapter 4 --- BDC additional and BMVR}
\emph{Folder:} \texttt{4 BDC additional and BMVR/document MAIN/}

\begin{enumerate}[leftmargin=2.2em,label=\textbf{4.\arabic*}]
\item \ERR\ \texttt{sections/09\_potential\_application\_in\_bmvr.tex}.
The central proposal $p=\langle X\rangle_{\rm QS}$ is dimensionally inconsistent:
a count substituted for a rate. See Part II \S6. The section admits at line 65 that
``The above formula is incorrect'' but attributes it to load-dependence of the burst
rate, which is a second-order issue.

\item \GAP\ Same file, line 60: ``(need to look up the formula for QS mean, but its
something like $\lim J(t)/(1-\Ihat(t))$. I think there is an important point around
whether $\mu=0$ or not.)'' --- The formula is $\lim J/\Ihat=a/(a-1)$
(Part I \S4.2); the $\mu$ point is real and is exactly that
$\langle X\rangle_{\rm QS}=V_\infty$ \emph{only} when $\mu=0$.

\item \GAP\ Same file: both candidate integrals are left unresolved. Neither works;
see Part II \S6 for why, and Part II \S7--\S8 for what replaces them.

\item \ERR\ Same file: $\Vfree(t)=N_0V(t)$ holds only for synchronously infected cells.
The section notices the asynchrony problem in the ``Towards a correction'' subsection but
does not act on it.

\item \GAP\ Same file: the closing subsection (``Iterating cycles'') describes the right
object in words --- a stationary regime of burst events with known time and size
distributions --- and stops. That object is the renewal kernel of Part II \S7.

\item \NOT\ Same file uses $R$ for the rupture state while Ch.3 uses $H$; and uses
$I$ for both the non-catastrophe probability and the infected-cell count in the same
displayed equations.

\item \TYP\ Same file: the macro \verb|\Vt| is defined in \texttt{main.tex} as
\verb|_{\text{total}} V|, which typesets a leading subscript with nothing to attach to.
Almost certainly not intended.

\item \GAP\ \texttt{sections/07\_burst\_size\_distribution.tex}: ``The current goal is to
verify this using Gillespie algorithm simulations.'' The result is correct and now has an
analytic proof plus two consistency checks (Part I \S4.4). Replace the note with the
proof; a simulation figure is then a bonus, not a dependency.

\item \TYP\ Same file: the antiderivative is first written with exponent $n$ and prefactor
$(n\beta)^{-1}$, then restated with $k$; and the same display appears twice.

\item \GAP\ \texttt{sections/06\_quasi\_stationary\_distribution.tex}: ``There is a
question of what form $F(t)$ takes, and how it differs between BDK and BDC.'' It is
determined by the section's own decomposition: $F=1-\Ihat=1-I+D$. Left open
unnecessarily.

\item \GAP\ Same file: the QSD is never identified. It is geometric$(1/a)$ on
$\{1,2,\dots\}$, for all $\mu$, and follows immediately from the section's own $p_n(t)$
(Part I \S4.2). This is the chapter's best unclaimed result.

\item \TYP\ Same file: the sentence ``Ignoring, for now, the case of $i=0$ \& $i=R$,
$\xt$ has reached absorption, $\Prob{\xt=i}=0$'' is garbled; the intended statement is
$\Prob{\xt=i\mid\text{fixation}}=0$ for $i\notin\{0,R\}$.

\item \TYP\ Same file, figure caption: two different curves are both described as
``yellow''.

\item \ERR/\GAP\ \texttt{sections/08\_further\_considerations.tex}: the ``curious
circularity'' is an artefact of using $\E[\Kb]$ for the conditional mean in the numerator
and the unconditional event in the denominator. Correct limit is $0$; the meaningful
conditional limit is $\langle X^2\rangle_{\rm QS}/\langle X\rangle_{\rm QS}$.
See Part I \S4.5. Also the ``(check)'' on $\E[\Kb]=V_{\rm inf}$ is now discharged.

\item \GAP\ \texttt{sections/03\_probability\_generating\_function\_pdes.tex}: in the
catastrophe variant $G(z,t)$ is a \emph{defective} PGF with $G(1,t)=I(t)\neq1$. This is
never stated and a reader will assume normalisation. The killing variant is proper.

\item \TYP\ Same file: ``The probability generating function (PDF), is defned'' --- both
the acronym and the spelling. Also ``derivd'', ``occurrs'', ``seperate'', ``labeled''
(inconsistent with British spelling elsewhere), ``killng''.

\item \STR\ \texttt{sections/01\_opening.tex} is empty but \verb|\input|.

\item \STR\ The folder contains a complete stale duplicate of Chapter~3 at
\texttt{4 BDC additional and BMVR/3 BDC core/}. Delete or archive.
\end{enumerate}

\section{Chapter 5 --- BDC Two Type}
\emph{Folder:} \texttt{5 BDC Two Type/Two\_Type\_Chapter/document MAIN/}

\begin{enumerate}[leftmargin=2.2em,label=\textbf{5.\arabic*}]
\item \GAP\ The two-type analogue of $D$ (extinction without catastrophe) and hence of
$\Ihat$ is absent. \S2 correctly notes that $S$ and $G$ ``are no-catastrophe probabilities
rather than population-survival probabilities'', which is exactly the $I$ versus $\Ihat$
distinction of Ch.3 --- but the two chapters never cross-reference it, and Ch.5 provides
only the $I$-analogue. Required for Part II. See Task~\ref{task:twotypeD}.

\item \GAP\ The killed second moments needed for a two-type release kernel are absent.
See Task~\ref{task:twotypemoments}.

\item \GAP\ GATE ($\delta_1=0<\delta_2$, $S'(0)=0$) is an eclipse phase. The chapter
describes GATE purely as a biological extreme and does not make the connection to the
within-host modelling literature, where the eclipse phase is a standing phenomenological
patch. This is a strong, cheap addition.

\item \STR\ \texttt{sections/06\_interpretation.tex} is a comment-only stub
(``Content folded into 04\_main\_result.tex'') but is still \verb|\input|.

\item \STR\ The proof of the main theorem is commented out in
\texttt{04\_main\_result.tex}. Deliberate or not, decide before submission.

\item \STR\ The chapter is written as a standalone paper (``The paper is organised
around the derivation''; ``A short conclusion and four technical appendices complete the
paper''). Either commit to it as a paper and give the thesis a wrapper, or convert the
self-references. Do not leave it ambiguous.

\item \NOT\ $\lambda_1,\lambda_2$ are birth rates here; $\lambda$ is the target-cell
influx in Ch.1 and the birth rate in Ch.3 \S3. $\mu_i$ are deaths, consistent.
$\delta_i$ are catastrophe rates, consistent with Ch.3/4 but colliding with Ch.1's
infected-cell death rate.

\item \NOT\ Absorbing state is $R$; Ch.3 uses $H$.

\item \GAP\ A figure caption embeds an external URL
(\fp{files.volvox.site/.../TT_bdc_genealogy_visualiser.html}). Confirm
permanence and that a link is acceptable to the examiners/journal, or move it to a
footnote with an archived copy.
\end{enumerate}

\section{Chapter 6 --- Bursting and budding}

\begin{enumerate}[leftmargin=2.2em,label=\textbf{6.\arabic*}]
\item \GAP\ \texttt{sections/01\_opening.tex} states that the BMVR is ``consistent with an
underlying stochastic model \dots{} each infected cell releases virions according to a
poisson process''. True, and the natural next sentence is that the resulting
\emph{offspring} distribution is geometric, not Poisson --- which is precisely the
comparator used in Part II \S\ref{sec:flooding}. Worth adding; it makes Ch.6 and the BMVR
extension continuous.

\item \GAP\ The intermediate release model needed for Task~8 (partial release; the
budding--bursting spectrum) is described qualitatively here. If it can be made explicit,
it is the bridge between this chapter and the BMVR extension.

\item \TYP\ ``ecclips phase'' $\to$ eclipse; ``wide speard'' $\to$ widespread;
``poisson'' $\to$ Poisson; ``non-cognisant''.
\end{enumerate}

\section{Cross-cutting}

\begin{enumerate}[leftmargin=2.2em,label=\textbf{X.\arabic*}]
\item \STR\ Each chapter is a standalone \texttt{article} with its own preamble, macro
block and \texttt{references.bib}. The macro blocks have already drifted apart
(e.g.\ \verb|\Vt|, \verb|\Ic| exist only in Ch.3/4). Unify into a shared
\texttt{preamble.tex} and a single \texttt{.bib} before the drift causes a real
inconsistency.

\item \STR\ Chapter cross-references throughout Ch.3 \S2 and Ch.1 \S6 use the old
numbering. Convert to \verb|\Chapref| / \verb|\cref| with real labels.

\item \NOT\ Adopt the canonical notation table of \S2 of these notes globally. The
$\beta$/$\lambda$, $\delta$/$\delta$, $I$/$I$ and $V$/$V$ collisions must be resolved
before Chapters 1 and 4 can be joined at all.

\item \STR\ Decide the identity of Chapter~5: thesis chapter, or paper with a thesis
wrapper. The same decision then applies to the BMVR extension of Part II.
\end{enumerate}

%%%%%%%%%%%%%%%%%%%%%%%%%%%%%%%%%%%%%%%%%%%%%%%%%%%%%%%%%%%%%%%%%%%%%%%%%%%%%%%
\section{Suggested paper outline}
%%%%%%%%%%%%%%%%%%%%%%%%%%%%%%%%%%%%%%%%%%%%%%%%%%%%%%%%%%%%%%%%%%%%%%%%%%%%%%%

Single paper, target \textit{Journal of Theoretical Biology} or
\textit{Bulletin of Mathematical Biology}.

\begin{enumerate}[leftmargin=1.8em,itemsep=0.2em]
\item \textbf{Introduction.} The BMVR's constant-release assumption; the single-cell
evidence against it; what is and is not new here (be explicit --- see Part II \S7.3).
\item \textbf{The intracellular process.} BDC; the closed forms of Part I; the QSD and
the burst-size law, with the identity between them.
\item \textbf{Renewal reformulation.} Equations \eqref{eq:Icell}--\eqref{eq:Vfree};
the kernels $\Ihat$ and $\delta K$.
\item \textbf{Exact reduction.} $p_{\rm eff}(r)$, $d_{\Icell,{\rm eff}}(r)$; the
characteristic equation; $R_0$ unchanged; $p$ depends on growth phase.
\item \textbf{Identifiability.} The negative result and the three rescuing observables.
\item \textbf{The flooding advantage.} Equation \eqref{eq:flood} and the condition
$1/\delta<1/\beta+1/\mu$. This is the headline.
\item \textbf{Application.} Single-cell release data; the caveat of Part II \S11.
\item \textbf{Discussion.} Limits of linear intracellular growth; the two-type/eclipse
extension as future work (or fold it in if Tasks \ref{task:twotypemoments} and
\ref{task:twotypeD} are completed in time).
\end{enumerate}

If the two-type tasks are completed, split into two papers: (i) the one-type theory as
above; (ii) the two-type kernel with a mechanistic eclipse phase, building on Chapter~5.

%%%%%%%%%%%%%%%%%%%%%%%%%%%%%%%%%%%%%%%%%%%%%%%%%%%%%%%%%%%%%%%%%%%%%%%%%%%%%%%
\appendix
\section{Verification record}
%%%%%%%%%%%%%%%%%%%%%%%%%%%%%%%%%%%%%%%%%%%%%%%%%%%%%%%%%%%%%%%%%%%%%%%%%%%%%%%

Every result in Parts I and II was checked symbolically and then again numerically at
$(\beta,\mu,\delta)\in\{(1,0.2,0.05),\,(1,0,0.1),\,(1,0.9,0.1),\,(1,0.5,\tfrac13)\}$,
and the stochastic claims were checked against an exact Gillespie simulation of the BDC
($3\times10^5$ realisations) and a direct simulation of the cell-level branching process
($6\times10^4$ realisations).

\begin{center}
\renewcommand{\arraystretch}{1.3}
\begin{tabular}{@{}lll@{}}
\toprule
Claim & Check & Outcome \\
\midrule
$ab=\mu/\beta$, $(a-1)(1-b)=\delta/\beta$ & symbolic & exact \\
Three forms of $V_\infty$ agree & symbolic + numeric & exact \\
$V(t)$ simplified $=$ thesis form & numeric, 5 times & $<10^{-10}$ \\
$\Ihat=(a-b)^2w/[(\Bb+\Aa w)(aw-b)]$ $=I-D$ & numeric & machine precision \\
$\sum_{n\ge1}p_n(t)=\Ihat(t)$ & symbolic + numeric & exact \\
$\sum_k\Prob{\Kb=k}=1-b$ & symbolic + numeric & exact \\
$\sum_k k\Prob{\Kb=k}=V_\infty$ & symbolic + numeric & exact \\
QSD $=$ burst-size law & symbolic & identical pmfs \\
$\langle X\rangle_{\rm QS}=a/(a-1)$ & numeric limit, $t=40$ & exact to $10^{-6}$ \\
$\int_0^\infty\Ihat=\beta^{-1}\log\frac{a}{a-1}$ & quadrature & $<2\times10^{-3}$ \\
$\E[\wt^2]$ corrected, eq.~\eqref{eq:EW2} & vs.\ exact geometric ($\mu=0$) & exact,
4 values of $\delta$ \\
$z_{\rm ext}^{\rm burst}$ formula & branching simulation & within $10^{-3}$ \\
Flooding identity \eqref{eq:flood} & numeric, 6 parameter sets & exact \\
\bottomrule
\end{tabular}
\end{center}

\subsection*{Simulation against theory, $(\beta,\mu,\delta)=(1,0.2,0.05)$}

\begin{center}
\renewcommand{\arraystretch}{1.3}
\begin{tabular}{@{}lrrl@{}}
\toprule
Quantity & Simulation & Theory & \\
\midrule
$\Prob{\text{clears, no burst}}$ & $0.18962$ & $0.18839$ $(=b)$ & \\
$I(1)$ & $0.92789$ & $0.92730$ & \\
$I(3)$ & $0.61773$ & $0.61597$ & \\
$I(8)$ & $0.20016$ & $0.19891$ & \\
$\E[W_\infty]$ & $13.957$ & $13.986$ $(=V_\infty)$ & \\
$\E[W_\infty^2]$ & $\mathbf{465.6}$ & $\mathbf{468.0}$ & \textbf{corrected
eq.~\eqref{eq:EW2}} \\
 & & $428.1$ & thesis eq.~(meanEY2) --- \emph{rejected} \\
$\E[\Kb\mid\text{burst}]$ & $17.223$ & $17.232$ $(=a/(a-1))$ & \\
$\Prob{\Kb=1\mid\text{burst}}$ & $0.05753$ & $0.05803$ & geometric$(1/a)$ \\
$\Prob{\Kb=5\mid\text{burst}}$ & $0.04615$ & $0.04569$ & \\
$\Prob{\Kb=25\mid\text{burst}}$ & $0.01379$ & $0.01382$ & \\
\bottomrule
\end{tabular}
\end{center}

The $\E[W_\infty^2]$ row is the decisive one: the corrected formula sits within
Monte-Carlo error of the simulation, the thesis formula is $8.5\%$ low --- far outside it.

\subsection*{Extinction probabilities, branching simulation}

\begin{center}
\renewcommand{\arraystretch}{1.3}
\begin{tabular}{@{}llrrr@{}}
\toprule
$(\beta,\mu,\delta)$ & $q$ & formula & simulation & budding \\
\midrule
$(1,0,0.1)$ & $0.2$ & $0.4000$ & $0.3986$ & $0.4545$ \\
$(1,0.9,0.1)$ & $0.9$ & $0.9351$ & $0.9351$ & $0.8442$ \\
$(1,0.5,\tfrac13)$ & $0.6$ & $0.8333$ & $0.8328$ & $0.8333$ \\
\bottomrule
\end{tabular}
\end{center}

The three rows are, respectively: bursting wins ($L>1$), bursting loses ($L<1$), and the
exact tie at $L=1$ predicted by \eqref{eq:flood}.

\subsection*{Numerical caution}

$\Ihat$ must be evaluated from the closed form \eqref{eq:Ihat}, not as $I-D$. Both tend
to $b$ while their difference tends to $0$, so the subtraction loses all significant
figures for $\theta t\gtrsim35$ --- at $(\beta,\mu,\delta)=(1,0.5,\tfrac13)$, $t=40$, the
subtraction returns $5\times10^{-4}$ where the true ratio $J/\Ihat$ is $3$. Anything
downstream that uses $\Ihat$ at large $t$ (notably $\int\Ihat$ and $\Lap{\Ihat}$) must use
\eqref{eq:Ihat}.

\end{document}
